\section{Chapter~\ref{sec:dofa}. Learning with \nt{DOPA}}

The mechanisms of the synapse (packets, coincidence detector, calcium, plasticity windows) and the behavior of dopamine as a prediction error are measured directly; that the rules lead to the gradient, and what broadcasting costs, are theorems; the outcome of learning to choose is computed with the book's model; the circuit that computes the error is a hypothesis.

{\footnotesize
\setlength{\tabcolsep}{3pt}\renewcommand{\arraystretch}{1.15}
\begin{longtable}{>{\raggedright}p{0.5cm}>{\raggedright}p{4.0cm}>{\raggedright}p{5.6cm}>{\raggedright}p{2.3cm}>{\raggedright\arraybackslash}p{1.7cm}}
\toprule
No. & Claim & Experiment and what was measured & Source & Status \\
\midrule\endfirsthead
\toprule
No. & Claim & Experiment and what was measured & Source & Status \\
\midrule\endhead
1 & Backpropagation of error in the form used in artificial networks has not been found in the brain & There is no direct experiment: this is a claim of absence. Reviews analyze what the algorithm requires (feedback connections mirroring the forward ones, a separate phase) and which biological approximations to it have been proposed. & \cite{Lillicrap2020, WhittingtonBogacz2019} & hypothesis \\
2 & The weight factors into the number of release sites, the release probability and the response to one packet: $w=N_s\,p\,A_1$~\eqref{eq:quantal} & Frog neuromuscular junction with reduced release: the recipient's response fluctuates in steps that are multiples of the spontaneous miniature response to one packet; the number of packets per spike obeys the Poisson law. & \cite{delCastillo1954} & measured \\
3 & The recipient's receptor is a coincidence detector; calcium triggers the weight change & Patch clamp of mouse neurons: Mg$^{2+}$ ions block the glutamate-opened channel at the resting potential and leave it upon depolarization. Hippocampal slices: a calcium chelator in the recipient blocks long-term potentiation, while light-triggered release of calcium inside the cell by itself strengthens transmission. & \cite{Nowak1984, Malenka1988calcium} & measured \\
4 & Either side carries out the decision: $A_1$ grows in the recipient, $p$ changes in the sender & Glutamate uncaged by light over a single spine causes rapid and lasting growth of the spine and of the current through its receptors; potentiation is accompanied by insertion of AMPA receptors. Depolarizing the recipient temporarily suppresses release by the sender; the effect is carried by endocannabinoids traveling back across the cleft (shown at \nt{GABA} inputs). Short-term depression depends on the sender's release probability. & \cite{Matsuzaki2004, Malinow2002, WilsonNicoll2001, Tsodyks1997} & measured \\
5 & A synapse strengthens when the sender and the recipient are active together~\eqref{eq:hebb} & Anesthetized rabbit: after brief trains of spikes along the hippocampal input pathway, the response is enhanced for $30$~min to $10$~h. In a hippocampal slice, the recipient's spikes strengthen only the synapses that are active at the same moment. & \cite{Bliss1973LTP, Kelso1986hebb} & measured \\
6 & Rule~\eqref{eq:oja} finds the principal eigenvector of the input correlations (Proposition~\ref{prop:oja}) & A theorem; the proof is in the chapter. There is no experiment in which a neuron has learned the principal component of its inputs; the mechanism that bounds weight growth in the synapse has not been established. & \cite{Oja1982} & theory \\
7 & Timing-based Hebb: a window of about $\pm20\ms$ (Fig.~\ref{fig:stdp}); chains grow from repeated sequences & Pairs of hippocampal neurons in culture: a recipient spike within $20\ms$ after a sender spike gives lasting strengthening, within $20\ms$ before, weakening; both depend on NMDA receptors. The ratio $A_-=0.6A_+$ in the figure is illustrative. That chains grow this way in a network has not been measured directly. & \cite{BiPoo1998} & measured \\
8 & Eligibility trace~\eqref{eq:threefactor}: dopamine arriving $1$--$2\,\text{s}$ after joint activity strengthens the connection, later dopamine does not (Proposition~\ref{prop:threefactor}) & Striatal slices, separate optogenetic stimulation of glutamate and dopamine inputs: the spine grows only if dopamine arrives $0.3$--$2\,\text{s}$ after the glutamate input; the window is set by the fast rise and decay of cAMP. In the cortex, pairing leaves separate traces for strengthening and weakening, which monoamine receptors convert into long-term changes within a window of order seconds. & \cite{Yagishita2014, He2015eligibility} & measured \\
9 & Windows lasting seconds also occur without dopamine: the third signal is strong excitation of the recipient & Awake mouse, area CA1: a single plateau event in the recipient strengthens inputs that arrived seconds before and after it, and creates a place field in a single pass. & \cite{Bittner2017} & measured \\
10 & The mean step of the three-factor rule follows the reward gradient~\eqref{eq:perturbation} (Proposition~\ref{prop:gradient}) & The gradient theorem for learning from random perturbations; derived in the chapter for small noise. & \cite{Williams1992} & theory \\
11 & Noise is search: the network learns from its own random deviations & Juvenile zebra finches: inactivating the output nucleus of the basal ganglia circuit abruptly and reversibly removes song variability. Adult zebra finches: noise is played on renditions of a syllable with pitch on one side of the mean, and the birds quickly shift the pitch the other way, learning from their own variability. & \cite{Olveczky2005, TumerBrainard2007} & measured \\
12 & A two-way choice is learned with $\tau_e=1\,\text{s}$ and $\delta=R-\bar R$; it is not learned with $\tau_e=50\ms$; without subtraction the weights grow without bound, and at a high learning rate the network can lock in the worse choice & \texttt{models/dofa\_learning.py}, $\eta=0.05$, $500$ trials, $6$ seeds. $\tau_e=1\,\text{s}$: fraction of $A$ choices $0.65$--$0.79\to0.96$--$1.00$, weights $0.85$--$1.33$. $\tau_e=50\ms$: $0.38$--$0.55$, weights unchanged. $\delta=R$: winner's weight $860$--$1190$. $\delta=R$, $\eta=0.5$, $3$ seeds: one locked onto $B$ ($0.04\to0$). The script prints all four cases; the rerun matches the chapter. & \texttt{models/\allowbreak dofa\_\allowbreak learning.py} & model \\
13 & The learning time from one shared signal grows in proportion to the number of noisy neurons (Proposition~\ref{prop:broadcastcost}) & Learning curves of a linear network: node perturbation is slower than the exact gradient by a factor equal to the number of output neurons, and weight perturbation is slower by a further factor equal to the number of inputs. & \cite{Werfel2005} & theory \\
14 & \nt{DOPA} outputs go primarily to action-selection regions & Complete axon reconstruction of single nigrostriatal \nt{DOPA} neurons in the rat: the branches of one cell cover $0.45$--$5.7\,\%$ of the striatal volume ($2.7\,\%$ on average), with almost no branches outside the striatum. & \cite{Matsuda2009} & measured \\
15 & The cortex learns to predict its input, and the error is computed on the spot & Review: sensory cortex shows responses to mismatch between expected and actual input. That this is a general learning rule of the cortex has not been proven. & \cite{Keller2018} & hypothesis \\
16 & Dopamine behaves like the error~\eqref{eq:ctd}: burst, transfer to the cue, dip (Fig.~\ref{fig:ctdcases}) & Monkey \nt{DOPA} neurons: a burst to unexpected juice; after learning, a burst to the cue and no response to the promised juice; a dip when the juice is omitted. The continuous form~\eqref{eq:ctd} is theory. & \cite{Schultz1997, Doya2000} & measured \\
17 & A circuit that computes $\dd V/\dd t$ and subtracts the expectation & Mice, recording and optogenetics in the ventral tegmental area: \nt{DOPA} neurons subtract the expectation from the reward; neighboring \nt{GABA} neurons inhibit them when reward is expected, and exciting or inhibiting these \nt{GABA} neurons changes the error. How the derivative is computed has not been shown. & \cite{Eshel2015arithmetic} & hypothesis \\
18 & The \nt{DOPA} neuron is a pacemaker; dips are shallower than bursts & In a slice, \nt{DOPA} neurons fire spontaneously and very regularly, riding on a slow pacemaker potential. In the monkey, the rate follows a positive error quantitatively but conveys a negative error only weakly. & \cite{GraceOnn1989, Bayer2005} & measured \\
19 & Actor and critic (Fig.~\ref{fig:actorcritic}) & The algorithm comes from reinforcement-learning theory. In humans (fMRI), the ventral striatum behaves like a critic both with and without choice, while the dorsal striatum does so only when actions must be chosen. & \cite{SuttonBarto2018, ODoherty2004actor} & hypothesis \\
20 & The sign of $\delta$ in the rule is flipped in neurons with the second receptor family & Striatal slices: in neurons with D1 receptors, pairing produces strengthening that requires D1; in neurons with D2, weakening that requires D2. & \cite{Shen2008} & measured \\
\bottomrule
\end{longtable}}
